\pagebreak
\thispagestyle{plain}

\begin{flushleft}
    \textbf{Elektrotehnički fakultet, Univerzitet u Sarajevu}\\
    \textbf{Odsjek za Računarstvo i informatiku}\\
    \textbf{Red. prof. dr. Haris Šupić, dipl. ing. el.}\\
\end{flushleft}

\begin{center}
    \vspace{2cm}
    {\Large Postavka zadatka završnog rada I ciklusa:}\\[0.2cm]
    {\large \textbf{Vektorske baze podataka:\\slučajevi upotrebe, algoritmi i ključne karakteristike}}
    \vspace{1cm}
\end{center}

\subsection*{Cilj:}

Cilj rada je istražiti teorijske osnove i praktične primjere primjene vektorskih baza podataka, te opis popularnih alata i sistema za upravljanje vektorskim bazama podataka.

\subsection*{Opis:}

Tema završnog rada obrađuje vektorske baze podataka s naglaskom na njihovu primjenu u savremenim informacionim sistemima. U radu je potrebno opisati vektorsku reprezentaciju podataka. Posebnu pažnju treba posvetiti načinu indeksiranja i pretraživanja vektora unutar baza podataka, te razlikama u odnosu na klasične relacijske baze. U teorijskom dijelu prikazuju se osnovni algoritmi za pretraživanje po sličnosti, s naglaskom na metode \textit{nearest neighbor} i \textit{approximate nearest neighbor}. Nadalje, u praktičnom dijelu rada daje se pregled popularnih alata i sistema za upravljanje vektorskim bazama kao što su Weaviate, FAISS i Milvus, uključujući i opis konkretnih primjena vektorskih baza podataka. Mogući primjeri mogu uključivati npr. sisteme za preporučivanje, pretraživanje slika, \textit{chatbot}-ove, sisteme za analizu dokumenata, ili neke druge samostalno odabrane primjere.

\subsection*{Očekivani rezultati:}

\begin{itemize}
    \item Opis teorijske podloge vektorskih baza podataka, uključujući njihove prednosti i razlike u odnosu na tradicionalne baze podataka.
    \item Opis osnovnih algoritama za pretraživanje po sličnosti, s naglaskom na metode \textit{nearest neighbor} i \textit{approximate nearest neighbor}.
    \item Prikaz primjera primjene vektorskih baza podataka.
\end{itemize}

\newpage

\subsection*{Polazna literatura:}

\begin{itemize}
    \item[] [1] A. Gionis, P. Indyk, and R. Motwani, "Similarity Search in High Dimensions via Hashing," in \textit{Proceedings of the 25th International Conference on Very Large Data Bases}, September 1999, pp. 518-529.
    \item[] [2] N. Ukey, Z. Yang, B. Li, G. Zhang, Y. Hu, and W. Zhang, "Survey on Exact kNN Queries over High-Dimensional Data Space," \textit{Sensors}, vol. 23, no. 2, p. 629, January 2023.
    \item[] [3] A. Guttman, "R-Trees: A Dynamic Index Structure for Spatial Searching," \textit{ACM SIGMOD Record}, vol. 14, no. 2, pp. 47-57, June 1984.
\end{itemize}

\begin{center}
    \vspace{2cm}
    \makebox[8cm]{\hrulefill}\\
    Red. prof. dr. Haris Šupić, dipl. ing. el.\\
\end{center}

\clearpage